\documentclass[12pt,a4paper]{article}

\usepackage[a4paper,margin=1in]{geometry}
\usepackage[T1]{fontenc}
\usepackage{lmodern}
\usepackage{xcolor}
\usepackage{graphicx}
\usepackage{eso-pic}
\usepackage{listings}
\usepackage{fancyhdr}
\usepackage{enumitem}

\newcommand{\addwatermark}{%
    \AddToShipoutPictureBG{%
        \AtPageLowerRight{%
            \raisebox{1.2cm}{%
                \makebox[0pt][r]{%
                    \rotatebox{45}{%
                        \textcolor{gray!20}{\fontsize{12}{14}\selectfont
                        Varun B Karmikanda\ --\ 261100610029}%
                    }%
                }%
            }%
        }%
        \AtPageUpperLeft{%
            \raisebox{-1.2cm}{%
                \hspace{1.5cm}\textcolor{white}{Varun B Karmikanda\ --\ 261100610029}%
            }%
        }%
    }%
}
\addwatermark

\definecolor{codegray}{gray}{0.95}
\lstset{
    language=C,
    basicstyle=\ttfamily\footnotesize,
    backgroundcolor=\color{codegray},
    frame=single,
    breaklines=true,
    showstringspaces=false,
    numbers=left,
    numberstyle=\tiny\color{gray},
    numbersep=8pt
}

\pagestyle{fancy}
\fancyhf{}
\lhead{RTOS LAB}
\rhead{Assignment 1}
\cfoot{\thepage}

\newcommand{\outputimage}[1]{%
    \begin{center}
    \includegraphics[width=0.9\textwidth,keepaspectratio]{images/#1}
    \end{center}
}

\begin{document}

\begin{titlepage}
    \centering
    \vspace*{2cm}
    {\Large\textbf{REAL TIME OPERATING SYSTEM}}\\[0.7cm]
    {\Large\textbf{RTOS LAB}}\\[1.2cm]
    {\Huge\textbf{Assignment 1}}\\[2cm]
    \begin{tabular}{r p{7.5cm}}
    \Large{Name:}     & {\Large\textbf{Varun B Karmikanda}} \\[0.4cm]
    \Large{USN:}      & {\Large{261100610029}} \\[0.4cm]
    \textbf{Program:} & ME Computer Science \& \\[0.1cm]
                      & Engineering
\end{tabular}
    \vfill
    {\large \today}
\end{titlepage}

\section*{1. LED Blink}
\textbf{Question:} Blink the onboard LED of the STM32 Blue Pill using
bare-metal programming.

\textbf{Code:}
\begin{lstlisting}
#include <stdint.h>
#include "stm32f103_regs.h"

static volatile uint32_t g_msticks;

void SysTick_Handler(void)
{
    g_msticks++;
}

static void delay_ms(uint32_t ms)
{
    uint32_t target = g_msticks + ms;
    while (g_msticks < target) {
    }
}

int main(void)
{
    RCC->APB2ENR |= RCC_APB2ENR_IOPCEN;

    GPIOC->CRH &= ~(0xFU << 20);
    GPIOC->CRH |= (GPIO_CNF_MODE_OUTPUT_PP_2MHZ << 20);

    SYSTICK->LOAD = (8000000U / 1000U) - 1U;
    SYSTICK->VAL = 0;
    SYSTICK->CTRL = SYSTICK_CTRL_CLKSOURCE |
                    SYSTICK_CTRL_TICKINT |
                    SYSTICK_CTRL_ENABLE;

    for (;;) {
        GPIOC->ODR ^= (1U << 13);
        delay_ms(1000);
    }
}
\end{lstlisting}

\textbf{Output:} The onboard PC13 LED blinks at a one-second interval.

\newpage
\section*{2. Three Thread Blink}
\textbf{Question:} Create three independent FreeRTOS tasks, each blinking
its own LED at a fixed interval:
\begin{itemize}[noitemsep]
    \item PC13: 1 second
    \item PA0: 2 seconds
    \item PA1: 3 seconds
\end{itemize}

\textbf{Code:}
\begin{lstlisting}
#include <stdint.h>
#include "stm32f103_regs.h"
#include "FreeRTOS.h"
#include "task.h"

typedef struct {
    GPIO_TypeDef *port;
    uint32_t pin;
    uint32_t period_ms;
} led_task_args_t;

static const led_task_args_t led_a = { GPIOC, 13, 1000 };
static const led_task_args_t led_b = { GPIOA, 0, 2000 };
static const led_task_args_t led_c = { GPIOA, 1, 3000 };
static const led_task_args_t led_d = { GPIOA, 2, 4000 };
static const led_task_args_t led_e = { GPIOA, 3, 5000 };

static void gpio_output_init(GPIO_TypeDef *port, uint32_t pin)
{
    volatile uint32_t *cr = (pin < 8) ? &port->CRL : &port->CRH;
    uint32_t shift = (pin % 8U) * 4U;
    *cr &= ~(0xFU << shift);
    *cr |= (GPIO_CNF_MODE_OUTPUT_PP_2MHZ << shift);
}

static void led_task(void *pvParameters)
{
    const led_task_args_t *args =
        (const led_task_args_t *)pvParameters;
    for (;;) {
        args->port->ODR ^= (1U << args->pin);
        vTaskDelay(pdMS_TO_TICKS(args->period_ms));
    }
}

int main(void)
{
    RCC->APB2ENR |= RCC_APB2ENR_IOPAEN | RCC_APB2ENR_IOPCEN;
    gpio_output_init(led_a.port, led_a.pin);
    gpio_output_init(led_b.port, led_b.pin);
    gpio_output_init(led_c.port, led_c.pin);
    gpio_output_init(led_d.port, led_d.pin);
    gpio_output_init(led_e.port, led_e.pin);

    xTaskCreate(led_task, "led_1s", configMINIMAL_STACK_SIZE,
                (void *)&led_a, tskIDLE_PRIORITY + 1, NULL);
    xTaskCreate(led_task, "led_2s", configMINIMAL_STACK_SIZE,
                (void *)&led_b, tskIDLE_PRIORITY + 1, NULL);
    xTaskCreate(led_task, "led_3s", configMINIMAL_STACK_SIZE,
                (void *)&led_c, tskIDLE_PRIORITY + 1, NULL);
    xTaskCreate(led_task, "led_4s", configMINIMAL_STACK_SIZE,
                (void *)&led_d, tskIDLE_PRIORITY + 1, NULL);
    xTaskCreate(led_task, "led_5s", configMINIMAL_STACK_SIZE,
                (void *)&led_e, tskIDLE_PRIORITY + 1, NULL);

    vTaskStartScheduler();
    for (;;) {
    }
}

void vApplicationStackOverflowHook(TaskHandle_t xTask, char *pcTaskName)
{
    (void)xTask;
    (void)pcTaskName;
    for (;;) {
    }
}
\end{lstlisting}

\textbf{Output:} The five configured task intervals are observed on the
connected LEDs.

\newpage
\section*{3. Three Thread Priority Switch}
\textbf{Question:} Create independent LED tasks with different priorities and
observe priority-based scheduling.

\textbf{Code:}
\begin{lstlisting}
#include <stdint.h>
#include "stm32f103_regs.h"
#include "FreeRTOS.h"
#include "task.h"

static volatile uint32_t g_msticks;

typedef struct {
    GPIO_TypeDef *port;
    uint32_t pin;
    uint32_t period_ms;
} led_task_args_t;

void vApplicationTickHook(void)
{
    g_msticks++;
}

static const led_task_args_t led_a = { GPIOA, 3, 500 };
static const led_task_args_t led_b = { GPIOA, 0, 450 };
static const led_task_args_t led_c = { GPIOA, 1, 400 };
static const led_task_args_t led_d = { GPIOA, 2, 350 };

static void delay_ms(uint32_t ms)
{
    uint32_t target = g_msticks + ms;
    while (g_msticks < target) {
    }
}

static void gpio_output_init(GPIO_TypeDef *port, uint32_t pin)
{
    volatile uint32_t *cr = (pin < 8) ? &port->CRL : &port->CRH;
    uint32_t shift = (pin % 8U) * 4U;
    *cr &= ~(0xFU << shift);
    *cr |= (GPIO_CNF_MODE_OUTPUT_PP_2MHZ << shift);
}

static void led_task(void *pvParameters)
{
    const led_task_args_t *args =
        (const led_task_args_t *)pvParameters;
    for (;;) {
        args->port->ODR ^= (1U << args->pin);
        delay_ms(args->period_ms);
    }
}

int main(void)
{
    RCC->APB2ENR |= RCC_APB2ENR_IOPAEN | RCC_APB2ENR_IOPCEN;
    gpio_output_init(led_a.port, led_a.pin);
    gpio_output_init(led_b.port, led_b.pin);
    gpio_output_init(led_c.port, led_c.pin);
    gpio_output_init(led_d.port, led_d.pin);

    xTaskCreate(led_task, "led_1s", configMINIMAL_STACK_SIZE,
                (void *)&led_a, tskIDLE_PRIORITY + 3, NULL);
    xTaskCreate(led_task, "led_2s", configMINIMAL_STACK_SIZE,
                (void *)&led_b, tskIDLE_PRIORITY + 2, NULL);
    xTaskCreate(led_task, "led_3s", configMINIMAL_STACK_SIZE,
                (void *)&led_c, tskIDLE_PRIORITY + 1, NULL);
    xTaskCreate(led_task, "led_4s", configMINIMAL_STACK_SIZE,
                (void *)&led_d, tskIDLE_PRIORITY + 0, NULL);

    vTaskStartScheduler();
    for (;;) {
    }
}

void vApplicationStackOverflowHook(TaskHandle_t xTask, char *pcTaskName)
{
    (void)xTask;
    (void)pcTaskName;
    for (;;) {
    }
}

void vApplicationIdleHook(void)
{
    __asm volatile ("wfi");
}
\end{lstlisting}

\textbf{Output:} The higher-priority task executes first whenever multiple
tasks are ready.

\newpage
\section*{4. Recursive and Addition Programs: GDB Trace and Debugging}
\textbf{Question:} Trace the recursive and addition programs using GDB output
and debug them.

\subsection*{4.1 Recursive Program}

\textbf{Code:}
\begin{lstlisting}
int rec_sum(int i)
{
    if (i <= 1)
        return i;
    else
        return i + rec_sum(i - 1);
}

int main(void)
{
    int sum = 0;
    sum = rec_sum(10);
    return 0;
}
\end{lstlisting}

\textbf{Output:} The calls proceed from \texttt{rec\_sum(10)} to
\texttt{rec\_sum(1)} and return through the stack. The final value is
\texttt{sum = 55}.
\outputimage{recursion-gdb-1.png}
\outputimage{recursion-gdb-2.png}
\outputimage{recursion-gdb-3.png}

\textbf{Trace summary:}
\[
1,\ 3,\ 6,\ 10,\ 15,\ 21,\ 28,\ 36,\ 45,\ 55
\]
The ARM disassembly shows the function prologue, stack-frame allocation,
argument storage, recursive branch, and return through register
\texttt{r0}.

\subsection*{4.2 Addition Program}

\textbf{Code:}
\begin{lstlisting}
int mySum(int a, int b)
{
    return a + b;
}

int main(void)
{
    int i = 5, j = 10, sum = 0;
    sum = mySum(i, j);
    return 0;
}
\end{lstlisting}

\textbf{Output:} The values \(i=5\) and \(j=10\) are passed in registers
\texttt{r0} and \texttt{r1}. The function returns \(15\) in \texttt{r0},
which is stored in \texttt{sum}.
\outputimage{sum-gdb.png}

\end{document}
